% ============================== WRAP-UP ==============================
\settab{WRAP-UP}{4}
\renewcommand{\chapribbontext}{Summary}
\chapter{The Whole Paper at a Glance}\label{ch:wrap}

\section{Four things to know by heart}

Almost every mark in this mock comes from one of these four statements. Write them from memory
before the real exam, until you can do it without hesitating.

\begin{defi}[Differentiability]
$f$ is differentiable at an interior point $P$ if
$\ f(P+H)=f(P)+\grad f(P)\cdot H+\norm H\,\varepsilon(H)$ with $\varepsilon(H)\to0$ as $H\to O$.
\end{defi}

\begin{thm}[The $C^1$ theorem]
If $f_x$ and $f_y$ exist and are continuous on an open set $U$, then $f$ is differentiable at every point of $U$.
\end{thm}

\begin{thm}[Chain rule]
If $X$ is differentiable at $t_0$ and $f$ is differentiable at $X(t_0)$, then
$\ \frac{d}{dt}f(X(t))\big|_{t_0}=\grad f(X(t_0))\cdot X'(t_0)$. For $u=\varphi(r(x,y))$:
$\ u_x=\varphi'(r)\,r_x$.
\end{thm}

\begin{thm}[Directional derivative]
If $f$ is differentiable at $P$ and $\norm u=1$, then $D_uf(P)=\grad f(P)\cdot u=\norm{\grad f(P)}\cos\theta$.
It is largest along $\grad f(P)$, smallest along $-\grad f(P)$, and zero along the level curve.
\end{thm}

\clearpage
\section{What the marker looks for}

The answers in this mock are short. The marks go to the justifications. Before you hand in, scan
your copy for these:

\begin{center}\small
\begin{tabularx}{\linewidth}{@{}>{\raggedright\arraybackslash}p{0.42\linewidth}Y@{}}\toprule
\textsf{\textbf{\color{tealDD}When you\ldots}} & \textsf{\textbf{\color{tealDD}\ldots make sure you wrote}}\\ \midrule
say a set is open or closed & the continuous function $g$ and the form $\{g>0\}$ or $\{g\ge0\}$; for ``not open'', one explicit point\\
square an equation & that both sides are $\ge0$\\
claim a function is differentiable & the open set, the partials, why they are continuous, ``by the $C^1$ theorem''\\
claim a function is \emph{not} differentiable & the candidate $\grad f(P)$, the quotient $\varepsilon(H)$, one path where $\varepsilon\not\to0$\\
use the chain rule & that $f$ is differentiable at the point where the curve is\\
compute a directional derivative & the unit vector; and if $f$ might not be differentiable, the limit definition instead of $\grad f\cdot u$\\
treat a function defined by cases & ``for $(x,y)\ne(0,0)$'' and ``at $(0,0)$'' as two separate paragraphs\\
conclude ``constant'' from a zero derivative & that the domain is an interval\\
give a limit at the origin & the cases (sign of a parameter) and the value when it exists\\ \bottomrule
\end{tabularx}
\end{center}

\section{Every answer in one table}

\begin{center}\small
\renewcommand{\arraystretch}{1.35}
\begin{tabularx}{\linewidth}{@{}l Y@{}}\toprule
\textsf{\textbf{\color{tealDD}Q}} & \textsf{\textbf{\color{tealDD}Answer}}\\ \midrule
\multicolumn{2}{@{}l}{\textsf{\textbf{\color{coralD}Problem 1}\quad $f=\sqrt{2y-x^2}$, $P=(1,1)$}}\\
1.1 & $D=\{y\ge\frac{x^2}2\}$, closed, not open\\
1.2 & $\varnothing$ if $c<0$; parabola $y=\frac{x^2}2$ if $c=0$; parabola $y=\frac{x^2+c^2}2$ if $c>0$\\
1.3 & $C^1$ on $\{2y-x^2>0\}$; $f(P)=1$, $\grad f(P)=(-1,1)$, $L(x,y)=1-x+y$\\
1.4 & plane $x-y+z=1$; normal line $(1+s,\,1-s,\,1+s)$, $s\in\R$\\
1.5 & $1-4t-t^2=5-(t+2)^2>0$; derivative $-2$\\
1.6 & $D_uf(P)=0$: $(1,1)$ is tangent to $C_1$, $\grad f(P)\perp(1,1)$\\
1.7 & increase $\frac1{\sqrt2}(-1,1)$ at rate $\sqrt2$; decrease $\frac1{\sqrt2}(1,-1)$ at rate $-\sqrt2$\\
\midrule
\multicolumn{2}{@{}l}{\textsf{\textbf{\color{coralD}Problem 2}\quad $g=\frac{xy^2}{x^2+y^2}$, $g(O)=0$}}\\
2.1 & $g_x=\frac{y^2(y^2-x^2)}{(x^2+y^2)^2}$, $g_y=\frac{2x^3y}{(x^2+y^2)^2}$; $g_x(O)=g_y(O)=0$\\
2.2 & $|g|\le|x|\to0$: continuous on $\R^2$\\
2.3 & $D_vg(O)=ab^2$\\
2.4 & not differentiable: $\varepsilon(t,t)=\frac1{2\sqrt2}\not\to0$\\
2.5 & $g_x(0,y)=1\not\to0$: $g_x$ not continuous at $O$, $g\notin C^1(\R^2)$\\
\midrule
\multicolumn{2}{@{}l}{\textsf{\textbf{\color{coralD}Problem 3}\quad $u=\varphi(r)$ on $\Omega$}}\\
3.1 & $u_x=\varphi'\frac xr$, $u_{xx}=\varphi''\frac{x^2}{r^2}+\varphi'\frac{y^2}{r^3}$ (and symmetric)\\
3.2 & $\Delta u=\varphi''+\frac1r\varphi'$\\
3.3 & $\Delta\ln r=0$; limit $-\infty$\\
3.4 & $(r\varphi')'=0$, so $r\varphi'=A$\\
3.5 & $u=A\ln\sqrt{x^2+y^2}+B$\\
3.6 & finite limit iff $A=0$; the limit is $B$\\ \bottomrule
\end{tabularx}
\end{center}

\clearpage
\section{Your two sittings}

Tick a box when you get the answer \emph{with its justification}, from blank paper, without looking.
The second column is for the sitting two days later. Anything still empty after the second sitting
is what you reread on the last evening.

\begin{center}\small
\renewcommand{\arraystretch}{1.5}
\newcommand{\bx}{\tikz\draw[tealD, line width=0.6pt, rounded corners=0.6pt] (0,0) rectangle (0.32,0.32);}
\begin{tabular}{@{}l c c @{\hspace{1.4cm}} l c c @{\hspace{1.4cm}} l c c@{}}\toprule
\multicolumn{3}{@{}l}{\textsf{\textbf{\color{coralD}Problem 1}}} & \multicolumn{3}{l}{\textsf{\textbf{\color{coralD}Problem 2}}} & \multicolumn{3}{l@{}}{\textsf{\textbf{\color{coralD}Problem 3}}}\\
 & \textsf{\scriptsize 1st} & \textsf{\scriptsize 2nd} & & \textsf{\scriptsize 1st} & \textsf{\scriptsize 2nd} & & \textsf{\scriptsize 1st} & \textsf{\scriptsize 2nd}\\ \midrule
1.1 & \bx & \bx & 2.1 & \bx & \bx & 3.1 & \bx & \bx\\
1.2 & \bx & \bx & 2.2 & \bx & \bx & 3.2 & \bx & \bx\\
1.3 & \bx & \bx & 2.3 & \bx & \bx & 3.3 & \bx & \bx\\
1.4 & \bx & \bx & 2.4 & \bx & \bx & 3.4 & \bx & \bx\\
1.5 & \bx & \bx & 2.5 & \bx & \bx & 3.5 & \bx & \bx\\
1.6 & \bx & \bx &     &     &     & 3.6 & \bx & \bx\\
1.7 & \bx & \bx &     &     &     &     &     &    \\ \bottomrule
\end{tabular}
\end{center}

\begin{remember}[title={The last evening}]
Do not learn anything new. Reread the four boxes at the top of this chapter, the ``lesson'' box of
any question you did not tick twice, and the ``What the marker looks for'' table. Then stop. You have
already sat this paper twice: in the exam room, it is the third time.
\end{remember}
